\section{Introdução}\label{sec:introducao}

Na Parte 1 foi desenvolvido o controle longitudinal do carrinho \cite{mozelli2025controlelongitudinal}. Para percorrer trajetórias curvilíneas, é necessário também um controle lateral, que atua no esterçamento das rodas dianteiras para manter o veículo sobre um caminho desejado. Neste trabalho foram implementados e comparados dois controladores laterais vistos em sala \cite{mozelli2025controlelateral}: um controlador \textit{on-off}, que comanda o esterçamento em três níveis ($-\delta_{on}$, $0$, $+\delta_{on}$) a partir de um erro de orientação, e o controlador Stanley \cite{thrun2006stanley,hoffmann2007stanley}, que calcula o esterçamento de forma contínua a partir dos erros de orientação e lateral no eixo dianteiro. Em todos os ensaios, a velocidade é mantida pelo PI longitudinal da Parte 1, sem alterações.

O desenvolvimento foi feito no CoppeliaSim \cite{coppeliasim}, na cena \texttt{simulador\_cones.ttt}, cuja pista tem uma linha amarela pintada no centro; os cones foram ignorados. Primeiro, os controladores foram testados com referências geométricas (reta, círculo e oito), usando a pose verdadeira do veículo, para avaliar as leis de controle sem erros de percepção. Em seguida, a referência passou a ser estimada somente pela câmera embarcada, por um algoritmo que segmenta a linha amarela e a projeta no chão, e o carrinho percorreu a volta completa da pista ($\approx246\,\mathrm{m}$) com diferentes velocidades, para determinar o quão rápido cada controlador completa o percurso sem capotar ou se perder. Por fim, o Stanley foi testado no carrinho real, seguindo uma faixa amarela no piso. Os códigos estão no repositório do projeto \cite{desgovernai2026}, e o arquivo \texttt{LEIA-ME.txt} descreve como reproduzir cada resultado.
